\documentclass[11pt,oneside]{article}

\usepackage[a4paper,margin=27mm,headheight=15pt]{geometry}
\usepackage[T1]{fontenc}
\usepackage[utf8]{inputenc}
\IfFileExists{libertinus.sty}{\usepackage{libertinus}}{\usepackage{lmodern}}
\usepackage{microtype}
\usepackage{amsmath,amssymb,amsthm,mathtools,mathrsfs}
\usepackage{bm}
\usepackage{booktabs,longtable,array,tabularx}
\usepackage{xcolor}
\usepackage{enumitem}
\usepackage{fancyhdr}
\usepackage{titlesec}
\usepackage{xurl}
\usepackage{hyperref}
\usepackage[nameinlink,noabbrev]{cleveref}

\definecolor{linkblue}{RGB}{25,75,135}
\definecolor{shadegray}{RGB}{246,247,249}
\definecolor{rulegray}{RGB}{195,201,208}
\definecolor{deepgray}{RGB}{70,75,82}
\hypersetup{
  colorlinks=true,
  linkcolor=linkblue,
  citecolor=linkblue,
  urlcolor=linkblue,
  pdftitle={Arithmetic Spectra of the Rvachev Sinc Product},
  pdfsubject={Thue--Morse lobe signs, exact lattice reproduction, complex dimensions, and q-binomial spectral layering},
  pdfkeywords={Fabius function, Rvachev up-function, Thue-Morse, sinc product, q-binomial, Strang-Fix, spectral zeta, Mellin transform}
}

\pagestyle{fancy}
\fancyhf{}
\fancyhead[L]{\small Fabius--Rvachev frontier report}
\fancyhead[R]{\small\nouppercase{\leftmark}}
\fancyfoot[C]{\thepage}
\renewcommand{\headrulewidth}{0.4pt}

\titleformat{\section}{\Large\bfseries}{\thesection}{0.7em}{}
  [\vspace{0.15ex}\color{rulegray}\titlerule]
\titleformat{\subsection}{\large\bfseries}{\thesubsection}{0.7em}{}
\titleformat{\subsubsection}{\normalsize\bfseries}{\thesubsubsection}{0.7em}{}
\titlespacing*{\section}{0pt}{2.8ex plus 1ex minus .2ex}{1.7ex plus .2ex}
\titlespacing*{\subsection}{0pt}{2.2ex plus .8ex minus .2ex}{1.0ex plus .2ex}

\setcounter{secnumdepth}{3}
\setcounter{tocdepth}{2}
\setlist[itemize]{leftmargin=2em,itemsep=.25em,topsep=.45em}
\setlist[enumerate]{leftmargin=2.3em,itemsep=.25em,topsep=.45em}
\setlength{\emergencystretch}{3em}
\allowdisplaybreaks[3]
\numberwithin{equation}{section}

\newtheorem{theorem}{Theorem}[section]
\newtheorem{proposition}[theorem]{Proposition}
\newtheorem{lemma}[theorem]{Lemma}
\newtheorem{corollary}[theorem]{Corollary}
\newtheorem{conjecture}[theorem]{Conjecture}
\theoremstyle{definition}
\newtheorem{definition}[theorem]{Definition}
\newtheorem{example}[theorem]{Example}
\newtheorem{problem}[theorem]{Open problem}
\theoremstyle{remark}
\newtheorem{remark}[theorem]{Remark}
\newtheorem{warning}[theorem]{Warning}

\newcommand{\R}{\mathbb R}
\newcommand{\C}{\mathbb C}
\newcommand{\N}{\mathbb N}
\newcommand{\Z}{\mathbb Z}
\newcommand{\E}{\mathbb E}
\newcommand{\Prob}{\mathbb P}
\newcommand{\dd}{\,\mathrm d}
\newcommand{\ii}{\mathrm i}
\newcommand{\e}{\mathrm e}
\newcommand{\Up}{\operatorname{up}}
\newcommand{\sinc}{\operatorname{sinc}}
\newcommand{\sgn}{\operatorname{sgn}}
\newcommand{\ord}{\operatorname{ord}}
\newcommand{\Res}{\operatorname*{Res}}
\newcommand{\supp}{\operatorname{supp}}
\newcommand{\Var}{\operatorname{Var}}
\newcommand{\bigO}{\mathcal O}
\newcommand{\smallo}{o}
\newcommand{\wt}{s_2}
\newcommand{\vtwo}{\nu_2}
\newcommand{\vb}{\nu_b}
\newcommand{\floor}[1]{\left\lfloor #1\right\rfloor}
\newcommand{\ceil}[1]{\left\lceil #1\right\rceil}
\newcommand{\abs}[1]{\left\lvert #1\right\rvert}
\newcommand{\norm}[1]{\left\lVert #1\right\rVert}
\newcommand{\qbinom}[3]{\genfrac{[}{]}{0pt}{}{#1}{#2}_{#3}}
\newcommand{\Corpus}{\textsc{Corpus}}
\newcommand{\Derived}{\textsc{Derived}}
\newcommand{\NewResult}{\textsc{New-to-snapshot}}
\newcommand{\Observed}{\textsc{Prior observation}}
\newcommand{\OpenStatus}{\textsc{Open}}
\newcommand{\repo}{\nolinkurl{Analysis/FabiusFunction/docs}}
\newcommand{\commit}{\nolinkurl{0b1db2f91321b6575be9a712dc1025e88af318db}}

\newenvironment{proofidea}
  {\par\smallskip\noindent\textit{Proof architecture.}\ }
  {\hfill$\square$\par\smallskip}
\newenvironment{boxedremark}{%
  \par\medskip\noindent\begin{center}\begin{minipage}{0.94\textwidth}%
  \color{deepgray}\hrule\vspace{0.65em}\small
}{%
  \vspace{0.65em}\hrule\end{minipage}\end{center}\medskip
}

\title{\bfseries Arithmetic Spectra of the Rvachev Sinc Product\\[0.35em]
\Large Thue--Morse lobe signs, exact lattice reproduction, complex dimensions,\\
and finite $q$-binomial spectral layering}
\author{Frontier report prepared for Vladimir Reshetnikov\\
\small after a recursive audit of the complete \texttt{ProveIt} Fabius-function TeX corpus}
\date{Repository snapshot \commit\\
\small 27 August 2026, 18:19:55 UTC}

\begin{document}
\maketitle

\begin{abstract}
This report starts from every \texttt{*.tex} document under \repo\ at the frozen
commit named above.  The corpus contains the principal Fabius--Rvachev exposition,
its historical revisions, a Lean walkthrough and glossary, exact dyadic and
$q$-calculus studies, repeated-integration and spline reports, endpoint asymptotics,
a large Thue--Morse formula atlas, a sequence of Fourier-decay audits, and primary
or transcribed papers.  Repeated revisions were read as one mathematical lineage;
the latest corrected statement was used whenever sibling drafts disagreed.

The principal synthesis is that the canonical-product zero multiplicities of the
Fourier transform
\[
 \Phi(z)=\widehat\Up(z)=\prod_{h\ge0}\sinc\!\left(\frac{\pi z}{2^h}\right)
       =\prod_{n\ge1}\left(1-\frac{z^2}{n^2}\right)^{1+\nu_2(n)}
\]
form an arithmetic spectrum.  Its multiplicity function
$a(n)=1+\nu_2(n)$ has Dirichlet series
$\zeta(s)/(1-2^{-s})$, exact counting law
$\sum_{n\le N}a(n)=2N-s_2(N)$, and a dense dyadic natural boundary in its
ordinary generating function.  These elementary facts organize several results
which do not appear as a theorem package in the pinned corpus.

First, the sign on the $N$th Fourier lobe is exactly the signed Thue--Morse bit
$(-1)^{s_2(N)}$.  This pattern had been observed before; here it is made a direct
corollary of the exact zero count and is converted into an absolutely convergent
Thue--Morse lobe decomposition of $\Up$ and all its derivatives.  Second, the
zero multiplicity gives an exact Strang--Fix hierarchy: on the lattice
$N^{-1}\Z$, polynomially weighted periodizations are constant through degree
$\nu_2(N)$, and the next degree fails.  A reciprocal-moment Appell correction then
reproduces every polynomial of that degree exactly.  Thus the approximation order
of the translate system is controlled not by the size of $N$, but by its
$2$-adic valuation.

Third, the zero spectrum has a spectral zeta function with ``complex dimensions''
$2\pi\ii k/\log2$.  Mellin inversion gives a superalgebraically accurate expansion
of $\log\Phi(\ii y)$ with the same Gamma--zeta Fourier coefficients that occur in
the repository's corrected endpoint asymptotics.  This supplies a new structural
explanation of that periodic fluctuation: it is the residue trace of the dyadic
zero spectrum.  Finite sinc products replace the geometric denominator by the
$q$-integer $[J]_{2^{-s}}$; the imaginary-axis poles are absent at every finite
stage and emerge through an explicit pole-condensation limit.  Gaussian
$q$-binomial coefficients appear as the algebra of independently selecting
spectral layers, while the common-zero statistics form a hierarchy
\[
 \sum_{n\ge1}\binom{1+\nu_2(n)}r n^{-s}
 =\zeta(s)\frac{2^{-(r-1)s}}{(1-2^{-s})^r}.
\]
A heat-trace expansion and a base-$b$ extension complete the report.

Claims marked \NewResult\ mean only that no counterpart was found in the frozen
TeX corpus after targeted collision searches.  They are proved here, but no claim
of worldwide historical priority is made.  In particular, Rvachev's 1977 paper
on approximation by $\Up$ and the broader Russian-language atomic-function
literature require a separate priority audit.
\end{abstract}

\begin{boxedremark}
\textbf{Status convention.}
\Corpus\ denotes an input already developed in the repository or cited primary
literature.  \Derived\ denotes a short consequence included to make the argument
self-contained.  \NewResult\ denotes a theorem not located in the pinned TeX
snapshot.  \Observed\ denotes a phenomenon recorded before but sharpened or
reorganized here.  None of the new-to-snapshot results is represented as already
formalized in Lean.
\end{boxedremark}

\tableofcontents
\newpage

\section{Audit boundary and research method}
\label{sec:audit}

\subsection{Why the snapshot is pinned}

The default branch changed while this audit was in progress.  To prevent a moving
target, all corpus statements in this report refer to commit \commit.  Its commit
time is 27 August 2026, 18:19:55 UTC.  GitHub's code index reported 54 TeX files
under \repo\ at audit time; because large generated and consolidated sources are
not always returned reliably by text search, the recursive Git tree, rather than
the index, was used as the authoritative inventory.

The documents fall into the following overlapping families.
\begin{longtable}{@{}p{0.24\textwidth}p{0.69\textwidth}@{}}
\toprule
Family & Mathematical role \\
\midrule
\endhead
Primary exposition & Definitions, probability model, functional equations,
canonical sinc product, exact dyadic values, derivative identities, partition of
unity, and the corrected global asymptotic framework. \\
Revision archive & Earlier expositions, supplements, glossaries, and walkthroughs;
useful for detecting abandoned normalizations and claims corrected in later waves. \\
Dyadic and $q$-calculus & Moment recurrences, Appell blocks, Gaussian-binomial
stencils, $q$-Pochhammer normalizations, exact dyadic arithmetic, Toeplitz limits,
and accelerated spline approximants. \\
Repeated integration & Finite nonuniform box splines, exact Thue--Morse truncated
powers, Peano kernels, and sharp $4^{-n}$ convergence laws. \\
Endpoint asymptotics & Correct Lambert phase, all-order saddle expansion, explicit
low-order coefficients, and the Gamma--zeta Fourier law for the log-periodic term. \\
Thue--Morse atlases & Digital, valuative, Pascal, automata, Mahler, finite-difference,
rarefaction, Mellin--Dirichlet, word-combinatorial, and spectral representations. \\
Fourier-decay audits & Successive corrections separating geometric, arithmetic,
RMS, and envelope exponents for the real-axis infinite sinc product. \\
Primary papers & Arias de Reyna's analytic and arithmetic papers, lacunary-product
work, Kronecker-sequence work, and OCR-corrected or transcribed source material. \\
\bottomrule
\end{longtable}

\subsection{Collision policy}

A formula is not called new merely because it admits a different notation.  The
following attractive directions were deliberately excluded from the novelty
claims.
\begin{itemize}
\item The all-order endpoint coefficients, including explicit $A_2$ and $A_3$,
      are already in the repository's small-argument asymptotics study.
\item Moment deconvolution and Richardson acceleration already occur in archived
      dyadic and convergence reports, including a half-base Gaussian-binomial
      filter across spline degrees.
\item Exact zero multiplicity $1+\nu_2(n)$ is already formalized in the Lean
      canonical-product module.
\item The sign pattern of the Fourier lobes was stated in Vladimir Reshetnikov's
      2018 Math StackExchange question, and the answer there essentially proves
      uniqueness of the critical point in every lobe from the logarithmic
      derivative.
\item Ordinary partition of unity is classical for $\Up$ and appears in the
      primary literature.  The new-to-snapshot statement below is the sharp
      higher weighted hierarchy and its exact $2$-adic stopping order.
\end{itemize}

The strategy was therefore to promote the zero multiplicity from an isolated
canonical-product fact to a spectrum and ask what its standard transforms know.
The answer simultaneously touches Dirichlet series, Mellin residues, automatic
sequences, shift-invariant approximation, and finite $q$-calculus.

\section{Normalization and established inputs}
\label{sec:normalization}

\subsection{The centered random series}

Let $U_1,U_2,\ldots$ be independent and uniform on $[-1,1]$, and set
\begin{equation}
 Y=\sum_{k=1}^{\infty}2^{-k}U_k.                    \label{eq:Y}
\end{equation}
Then $Y$ is supported on $[-1,1]$ and has the even density $\Up$.  Under the
Fourier convention
\begin{equation}
 \widehat f(\xi)=\int_{\R}f(x)\e^{-2\pi\ii x\xi}\dd x,      \label{eq:fourier-conv}
\end{equation}
independence gives
\begin{equation}
 \Phi(\xi):=\widehat\Up(\xi)
 =\prod_{h=0}^{\infty}\sinc\!\left(\frac{\pi\xi}{2^h}\right),
 \qquad \sinc z=\frac{\sin z}{z}.                          \label{eq:Phi-sinc}
\end{equation}
The convergence is locally uniform on $\C$, so $\Phi$ is an even entire function.
The density is $C^\infty$ and compactly supported; consequently $\Phi$ is rapidly
decreasing on the real axis, although its sharp pointwise behavior is highly
nonuniform.

The bounded Fabius variable is
\begin{equation}
 X=\frac{Y+1}{2}=\sum_{k=1}^{\infty}2^{-k}V_k,
 \qquad V_k\sim\operatorname{Unif}[0,1],                  \label{eq:X}
\end{equation}
with distribution function $F$.  This affine relation will be the bridge between
the imaginary-axis growth of $\Phi$ and the small-argument asymptotics of $F$.

\subsection{Canonical product and exact orders}

Euler's product for sine and an absolutely convergent rearrangement give
\begin{equation}
 \Phi(z)=\prod_{n=1}^{\infty}
 \left(1-\frac{z^2}{n^2}\right)^{a(n)},
 \qquad a(n)=1+\nu_2(n).                                  \label{eq:canonical}
\end{equation}
Thus
\begin{equation}
 \ord_{z=n}\Phi=a(n),\qquad n\in\Z\setminus\{0\}.          \label{eq:zero-order}
\end{equation}
This identity is an established corpus input.  Everything in
\cref{sec:spectrum,sec:weighted,sec:mellin,sec:qfinite} is driven by taking it
seriously as spectral data rather than merely as a description of zeros.

The first centered moments needed later are
\begin{equation}
 \mu_0=1,\qquad \mu_1=0,\qquad \mu_2=\frac19,
 \qquad \mu_3=0,\qquad \mu_4=\frac{19}{675},               \label{eq:moments}
\end{equation}
where $\mu_r=\int x^r\Up(x)\dd x$.  They follow either from the random series or
from the repository's moment recurrence.

\section{The arithmetic zero spectrum}
\label{sec:spectrum}

\subsection{Equivalent descriptions}

\begin{theorem}[Arithmetic spectrum package \NewResult]
\label{thm:spectrum-package}
For $n\ge1$, let $a(n)=1+\nu_2(n)$.  Then:
\begin{align}
 a(n)&=\sum_{h\ge0}\mathbf 1_{2^h\mid n},                  \label{eq:a-divisor}\\
 \sum_{n\ge1}\frac{a(n)}{n^s}
   &=\frac{\zeta(s)}{1-2^{-s}},\qquad \Re s>1,             \label{eq:D}\\
 A(N):=\sum_{n\le N}a(n)
   &=\sum_{h\ge0}\floor{\frac{N}{2^h}}
     =2N-s_2(N),                                           \label{eq:AN}\\
 \#\{n\le N:a(n)\ge r\}
   &=\floor{\frac{N}{2^{r-1}}},                            \label{eq:tail-count}\\
 \#\{n\le N:a(n)=r\}
   &=\floor{\frac{N}{2^{r-1}}}-
     \floor{\frac{N}{2^r}}.                               \label{eq:exact-count}
\end{align}
Moreover $a$ is multiplicative and satisfies
\begin{equation}
 a(2n)=a(n)+1,\qquad a(2n+1)=1.                            \label{eq:a-rec}
\end{equation}
Its limiting value distribution is geometric:
\begin{equation}
 \lim_{N\to\infty}\frac1N\#\{n\le N:a(n)=r\}=2^{-r},
 \qquad r\ge1,                                            \label{eq:geom}
\end{equation}
so the limiting mean and variance are both $2$.
\end{theorem}

\begin{proof}
The number of nonnegative $h$ for which $2^h\mid n$ is exactly
$1+\nu_2(n)$, proving \eqref{eq:a-divisor}.  Absolute convergence permits
summation by layers:
\[
 \sum_{n\ge1}\frac{a(n)}{n^s}
 =\sum_{h\ge0}\sum_{2^h\mid n}n^{-s}
 =\zeta(s)\sum_{h\ge0}2^{-hs}.
\]
The same layer count truncated at $N$ gives the first expression for $A(N)$.
Legendre's digit-sum identity
$\sum_{h\ge1}\floor{N/2^h}=N-s_2(N)$ gives the last one.
The inequalities $a(n)\ge r$ and $2^{r-1}\mid n$ are equivalent, giving
\eqref{eq:tail-count}; subtraction gives \eqref{eq:exact-count}.
Multiplicativity follows because only the $2$-part contributes: if $(m,n)=1$,
at most one is even, and $a(mn)=a(m)a(n)$.  The recursion and limiting law are
immediate.  The geometric mean is $\sum r2^{-r}=2$ and its variance is $2$.
\end{proof}

\begin{remark}
The average zero multiplicity is therefore $2$, while exceptional powers of two
have unbounded multiplicity.  This is the first indication that average real-axis
behavior and extremal dyadic behavior should not be conflated--precisely the
lesson reached independently by the repository's Fourier-decay audit sequence.
\end{remark}

\subsection{Ordinary generating function and natural boundary}

\begin{proposition}[Lambert series and a dyadic natural boundary \NewResult]
\label{prop:ogf}
For $|z|<1$,
\begin{equation}
 \mathcal A(z):=\sum_{n\ge1}a(n)z^n
 =\sum_{h\ge0}\frac{z^{2^h}}{1-z^{2^h}},                  \label{eq:A-ogf}
\end{equation}
with Mahler recursion
\begin{equation}
 \mathcal A(z)=\frac{z}{1-z}+\mathcal A(z^2).              \label{eq:A-mahler}
\end{equation}
Every root of unity of order a power of two is a radial singularity.  Since these
roots are dense on the unit circle, $|z|=1$ is a natural boundary.
\end{proposition}

\begin{proof}
Expand \eqref{eq:a-divisor} and sum geometric series.  Splitting off $h=0$
and shifting the remaining index gives \eqref{eq:A-mahler}.  If $\omega^{2^m}=1$
and $z=r\omega$ with $r\uparrow1$, the $h=m$ summand is
$r^{2^m}/(1-r^{2^m})\to+\infty$.  Thus analytic continuation across any arc is
impossible.
\end{proof}

\subsection{Spectral traces and the Taylor law}

\begin{corollary}[Even spectral traces \NewResult]
\label{cor:traces}
For every integer $m\ge1$,
\begin{equation}
 \sum_{n\ge1}\frac{a(n)}{n^{2m}}
 =\frac{\zeta(2m)}{1-4^{-m}},                              \label{eq:trace}
\end{equation}
and, for $|z|<1$,
\begin{equation}
 \log\Phi(z)
 =-\sum_{m=1}^{\infty}
 \frac{\zeta(2m)}{m(1-4^{-m})}z^{2m}.                    \label{eq:log-taylor}
\end{equation}
\end{corollary}

\begin{proof}
The first formula is \eqref{eq:D} at $s=2m$.  Expanding
$\log(1-z^2/n^2)$ in the canonical product and interchanging the absolutely
convergent sums gives the second.
\end{proof}

The logarithmic derivative is the corresponding resolvent trace:
\begin{equation}
 \frac{\Phi'(z)}{\Phi(z)}
 =\sum_{n\ge1}a(n)\left(\frac1{z-n}+\frac1{z+n}\right)
 =-2z\sum_{n\ge1}\frac{a(n)}{n^2-z^2}.                   \label{eq:logder}
\end{equation}
This identity will control the geometry of every real lobe.

\section{Thue--Morse signs from zero counting}
\label{sec:signs}

\subsection{The neighboring-bit identity}

Let
\begin{equation}
 \varepsilon_n=(-1)^{s_2(n)},\qquad n\ge0,                 \label{eq:eps}
\end{equation}
be the signed Thue--Morse sequence.

\begin{lemma}[Multiplicity parity is a Thue--Morse discrete derivative]
\label{lem:parity}
For every $n\ge1$,
\begin{equation}
 (-1)^{a(n)}=\varepsilon_{n-1}\varepsilon_n.              \label{eq:neighbor}
\end{equation}
\end{lemma}

\begin{proof}
Adding one to $n-1$ erases $\nu_2(n)$ trailing ones and creates one new one, so
$s_2(n)=s_2(n-1)+1-\nu_2(n)$.  Hence
$s_2(n-1)+s_2(n)\equiv1+\nu_2(n)=a(n)\pmod2$.
\end{proof}

\subsection{Exact lobe signs}

\begin{theorem}[Thue--Morse lobe-sign theorem \Observed]
\label{thm:lobe-sign}
For real $x\notin\Z$,
\begin{equation}
 \sgn\Phi(x)=(-1)^{s_2(\floor{|x|})}
             =\varepsilon_{\floor{|x|}}.                  \label{eq:lobe-sign}
\end{equation}
Equivalently, the signs of the consecutive positive-axis lobes of the infinite
sinc product are exactly the signed Thue--Morse word.
\end{theorem}

\begin{proof}
The product is positive at the origin.  Crossing the zero at $n$ reverses sign
exactly when $a(n)$ is odd.  Therefore, for $N<x<N+1$,
\[
 \sgn\Phi(x)=(-1)^{A(N)}.
\]
By \eqref{eq:AN}, $A(N)=2N-s_2(N)$, whose parity is $s_2(N)$.  Evenness handles
negative $x$.
\end{proof}

\begin{remark}
The pattern itself predates this report.  The point of
\cref{thm:lobe-sign} is the exact spectral explanation: Thue--Morse is not an
external decoration placed on the sinc product.  It is the parity integral of
its $2$-adic zero multiplicities, as \eqref{eq:neighbor} makes literal.
\end{remark}

For the $J$-factor truncation
\begin{equation}
 \Phi_J(x)=\prod_{h=0}^{J-1}\sinc\!\left(\frac{\pi x}{2^h}\right),             \label{eq:PhiJ}
\end{equation}
we likewise have, away from its zeros,
\begin{equation}
 \sgn\Phi_J(x)=(-1)^{\sum_{h=0}^{J-1}\floor{|x|/2^h}}.     \label{eq:finite-sign}
\end{equation}
Thus the infinite Thue--Morse sign law is already visible through finite digital
floor sums.

\subsection{A Thue--Morse lobe decomposition of the bump}

\begin{theorem}[Signed lobe inversion \NewResult]
\label{thm:lobe-inversion}
For every real $x$,
\begin{equation}
 \Up(x)=2\sum_{N=0}^{\infty}\varepsilon_N
 \int_N^{N+1}|\Phi(\xi)|\cos(2\pi x\xi)\dd\xi.             \label{eq:lobe-inversion}
\end{equation}
The series and every termwise differentiated series converge absolutely.  In
particular, for $m\ge0$,
\begin{equation}
 \Up^{(2m)}(0)=2(-1)^m(2\pi)^{2m}
 \sum_{N=0}^{\infty}\varepsilon_N
 \int_N^{N+1}\xi^{2m}|\Phi(\xi)|\dd\xi.                  \label{eq:lobe-moments}
\end{equation}
For $m=0$ this becomes the exact alternating-area identity
\begin{equation}
 1=2\sum_{N\ge0}\varepsilon_N\int_N^{N+1}|\Phi(\xi)|\dd\xi.                \label{eq:area}
\end{equation}
\end{theorem}

\begin{proof}
Fourier inversion and evenness give
$\Up(x)=2\int_0^\infty\Phi(\xi)\cos(2\pi x\xi)\dd\xi$.
Partition the half-line into unit lobes and apply \cref{thm:lobe-sign}.
Rapid decay of $\Phi$ permits arbitrary polynomial weights, proving absolute
convergence and termwise differentiation.
\end{proof}

This identity is a genuinely analytic bridge back from the sinc product to the
Thue--Morse sequence.  It differs from the better-known appearance of
Thue--Morse in truncated-power formulas for finite box splines: here the sequence
weights frequency lobes of the limiting $C^\infty$ bump itself.

\section{Laguerre--Polya geometry and real critical points}
\label{sec:LP}

\begin{theorem}[Real-rooted derivative hierarchy \Derived]
\label{thm:LP}
The entire function $\Phi$ belongs to the Laguerre--Polya class.  Consequently,
every derivative $\Phi^{(r)}$ has only real zeros.  More precisely:
\begin{enumerate}
\item the zero $n\in\Z\setminus\{0\}$ is inherited by $\Phi'$ with multiplicity
      $a(n)-1$;
\item in every interval $(n,n+1)$, $n\ge0$, $\Phi'$ has exactly one simple zero
      $c_n$;
\item for every $r\ge0$ and nonzero integer $m$,
      \begin{equation}
       \Phi^{(r)}(2^r m)=0.                                \label{eq:derivative-lattice}
      \end{equation}
\end{enumerate}
\end{theorem}

\begin{proof}
The canonical product is the locally uniform limit of real-rooted even
polynomials, because $\sum_n a(n)n^{-2}<\infty$.  Hence it is in the
Laguerre--Polya class, which is closed under differentiation.

An order-$a(n)$ zero loses one order under differentiation.  On $(n,n+1)$,
\eqref{eq:logder} is strictly decreasing, since
\[
 \left(\frac{\Phi'}{\Phi}\right)'(x)
 =-\sum_{m\ge1}a(m)\left(\frac1{(x-m)^2}+\frac1{(x+m)^2}\right)<0.
\]
It tends to $+\infty$ at $n+$ and to $-\infty$ at $(n+1)-$, so it has one
simple zero.  Finally, the zero at $2^rm$ has multiplicity at least $r+1$.
\end{proof}

\begin{corollary}[Signed unique extrema]
The unique critical value in the $n$th lobe has sign $\varepsilon_n$.  Thus the
critical-value sequence may be written
\begin{equation}
 \Phi(c_n)=\varepsilon_n M_n,\qquad M_n>0.                 \label{eq:critical-sign}
\end{equation}
\end{corollary}

The repository's Fourier-decay program studies the highly irregular envelope
$M_n$.  The arithmetic spectrum explains why a single ``typical'' exponent cannot
capture all subsequences: the inherited multiplicity at each endpoint varies
with $\nu_2(n)$, while the lobe sign varies with the integrated parity
$s_2(n)\bmod2$.

\section{Exact $2$-adic Strang--Fix hierarchy}
\label{sec:weighted}

This section contains the strongest new-to-snapshot theorem.  The ordinary
partition of unity is the degree-zero case.  The full hierarchy reveals an exact
and sharp dependence on the $2$-adic valuation of the lattice denominator.

\subsection{Weighted periodizations}

For $N\ge1$ and $r\ge0$, define
\begin{equation}
 \mathcal P_{N,r}(x)=
 \sum_{k\in\Z}\left(x-\frac{k}{N}\right)^r
 \Up\!\left(x-\frac{k}{N}\right).                         \label{eq:periodization}
\end{equation}
The sum is locally finite because $\Up$ has compact support.

\begin{theorem}[Sharp weighted partition hierarchy \NewResult]
\label{thm:weighted}
Let $q=\nu_2(N)$.  For every $0\le r\le q$ and every real $x$,
\begin{equation}
 \boxed{\mathcal P_{N,r}(x)=N\mu_r.}                       \label{eq:weighted-constant}
\end{equation}
The statement is sharp: $\mathcal P_{N,q+1}$ is not constant.  Equivalently, the
translate system generated by $\Up$ on the lattice $N^{-1}\Z$ has exact
Strang--Fix order
\begin{equation}
 \operatorname{SF}(N)=1+\nu_2(N).                          \label{eq:SF-order}
\end{equation}
\end{theorem}

\begin{proof}
Put $g_r(t)=t^r\Up(t)$.  Poisson summation for the $N^{-1}$-periodization gives
\begin{equation}
 \mathcal P_{N,r}(x)=N\sum_{m\in\Z}
 \widehat g_r(mN)\e^{2\pi\ii mNx}.                        \label{eq:poisson-weighted}
\end{equation}
Differentiating the Fourier transform,
\begin{equation}
 \widehat g_r(\xi)=\frac{\Phi^{(r)}(\xi)}{(-2\pi\ii)^r}.  \label{eq:moment-fourier}
\end{equation}
For $m\ne0$, the zero at $mN$ has order
\[
 a(mN)=1+\nu_2(mN)=1+q+\nu_2(m)\ge q+1.
\]
Hence $\Phi^{(r)}(mN)=0$ for $r\le q$.  Only $m=0$ remains in
\eqref{eq:poisson-weighted}, giving $N\widehat g_r(0)=N\mu_r$.
At $r=q+1$, the Fourier coefficient with $m=1$ is nonzero because the zero at
$N$ has exact order $q+1$.  Therefore the periodization cannot be constant.
\end{proof}

\begin{example}[The first three layers]
For every $N$,
\begin{equation}
 \sum_{k\in\Z}\Up(x-k/N)=N.                               \label{eq:PU}
\end{equation}
If $2\mid N$, then
\begin{equation}
 \sum_{k\in\Z}(x-k/N)\Up(x-k/N)=0.                        \label{eq:first-weight}
\end{equation}
If $4\mid N$, then
\begin{equation}
 \sum_{k\in\Z}(x-k/N)^2\Up(x-k/N)=\frac{N}{9}.           \label{eq:second-weight}
\end{equation}
The quadratic identity fails for $N\equiv2\pmod4$ and the linear identity fails
for odd $N$.
\end{example}

\begin{remark}[Why the result is arithmetically unusual]
Refining the lattice from $1/N$ to $1/(N+1)$ can reduce the exact reproduction
order dramatically.  Powers of two are optimal: the lattice $2^{-q}\Z$ has
order $q+1$, whereas every odd denominator has only the classical order one.
\end{remark}

\subsection{Exact polynomial reproduction}

Let
\begin{equation}
 M(z)=\int_{\R}\e^{zt}\Up(t)\dd t
     =\sum_{r\ge0}\mu_r\frac{z^r}{r!}.                    \label{eq:mgf}
\end{equation}
Since $\Up$ is even, $M(-z)=M(z)$.  Define reciprocal-moment Appell polynomials
$P_r$ by
\begin{equation}
 \frac{\e^{xz}}{M(z)}=\sum_{r=0}^{\infty}P_r(x)\frac{z^r}{r!}.              \label{eq:Pr}
\end{equation}

\begin{theorem}[Appell-corrected lattice reproduction \NewResult]
\label{thm:reproduction}
If $q=\nu_2(N)$, then for every polynomial degree $0\le r\le q$,
\begin{equation}
 \boxed{\frac1N\sum_{k\in\Z}P_r(k/N)\Up(x-k/N)=x^r.}      \label{eq:reproduce}
\end{equation}
Consequently every polynomial of degree at most $q$ is reproduced exactly by a
locally finite linear combination of $N^{-1}$-lattice translates of $\Up$.
\end{theorem}

\begin{proof}
Write $t_k=x-k/N$.  Up to degree $q$, \cref{thm:weighted} implies the formal
Taylor identity
\[
 \frac1N\sum_k\e^{-zt_k}\Up(t_k)=M(-z)+\bigO(z^{q+1})
 =M(z)+\bigO(z^{q+1}).
\]
Multiplying by $\e^{xz}/M(z)$ yields
\[
 \frac1N\sum_k\frac{\e^{(x-t_k)z}}{M(z)}\Up(t_k)
 =\e^{xz}+\bigO(z^{q+1}).
\]
Comparing coefficients gives \eqref{eq:reproduce}.
\end{proof}

The first correction polynomials are
\begin{align}
 P_0(x)&=1,\nonumber\\
 P_1(x)&=x,\nonumber\\
 P_2(x)&=x^2-\frac19,\nonumber\\
 P_3(x)&=x^3-\frac{x}{3},\nonumber\\
 P_4(x)&=x^4-\frac23x^2+\frac{31}{675}.                   \label{eq:P-first}
\end{align}
Thus, for example, the $1/16$ lattice reproduces quartics exactly.

\subsection{A formalization-friendly version}

The analytic proof above uses Poisson summation, but the theorem can be decomposed
for Lean into three reusable statements:
\begin{enumerate}
\item Fourier transform of $x^r\Up(x)$ equals
      $(-2\pi\ii)^{-r}\Phi^{(r)}$;
\item periodization on $N^{-1}\Z$ has Fourier coefficients sampled at $mN$;
\item exact zero order at $mN$ is $1+\nu_2(mN)$.
\end{enumerate}
No asymptotic estimate is needed.  Once a suitable Poisson theorem for compactly
supported smooth functions is available, the sharpness proof is a one-coefficient
argument.

\section{Spectral zeta function and complex dimensions}
\label{sec:zeta}

\begin{definition}[Zero spectral zeta function]
For $\Re s>1$, set
\begin{equation}
 Z_\Phi(s)=\sum_{n\ge1}\frac{\ord_n\Phi}{n^s}
          =\sum_{n\ge1}\frac{a(n)}{n^s}.                  \label{eq:Zphi}
\end{equation}
\end{definition}
By \eqref{eq:D},
\begin{equation}
 Z_\Phi(s)=\frac{\zeta(s)}{1-2^{-s}}.                     \label{eq:Zphi-form}
\end{equation}
It extends meromorphically to $\C$.  Besides the ordinary density pole at $s=1$,
it has the dyadic lattice of poles
\begin{equation}
 \chi_k=\frac{2\pi\ii k}{\log2},\qquad k\in\Z,            \label{eq:chi}
\end{equation}
except where an accidental zeta zero would cancel one.  At $k\ne0$,
\begin{equation}
 \Res_{s=\chi_k}Z_\Phi(s)=\frac{\zeta(\chi_k)}{\log2}.    \label{eq:resZ}
\end{equation}
We call the $\chi_k$ the complex dimensions of the dyadic zero spectrum.  This
terminology is descriptive: the underlying object is simply the meromorphic
Dirichlet series \eqref{eq:Zphi-form}.

\begin{remark}
The pole at $s=1$ has residue $2$, agreeing with the exact zero count
$A(N)=2N-s_2(N)$.  The imaginary lattice records the discrete scale invariance
which is invisible in the leading density but reappears as a periodic function
of $\log_2$ scale.
\end{remark}

\section{Mellin transform on the imaginary axis}
\label{sec:mellin}

\subsection{Exact transform}

Set
\begin{equation}
 H(y)=\log\Phi(\ii y)
     =\sum_{n\ge1}a(n)\log\left(1+\frac{y^2}{n^2}\right),
 \qquad y>0.                                               \label{eq:H}
\end{equation}
All terms are positive, so $H$ has none of the real-axis cancellation difficulties.

\begin{theorem}[Mellin factorization \NewResult]
\label{thm:Mellin}
For $1<\Re s<2$,
\begin{equation}
 \int_0^\infty H(y)y^{-s-1}\dd y
 =\frac{\pi\zeta(s)}{s\sin(\pi s/2)(1-2^{-s})}.            \label{eq:Mellin}
\end{equation}
Equivalently, with $S=2\pi y$,
\begin{equation}
 H(y)=\frac1{2\pi\ii}\int_{c-\ii\infty}^{c+\ii\infty}
 \frac{\Gamma(1-s)\zeta(1-s)}{s(1-2^{-s})}S^s\dd s,
 \qquad 1<c<2.                                             \label{eq:Mellin-inv}
\end{equation}
\end{theorem}

\begin{proof}
For $0<\Re s<2$, integration by parts and the beta integral give
\begin{equation}
 \int_0^\infty\log(1+t^2)t^{-s-1}\dd t
 =\frac{\pi}{s\sin(\pi s/2)}.                             \label{eq:kernel-Mellin}
\end{equation}
For $\Re s>1$, Tonelli's theorem and $y=nt$ factor the sum into this kernel and
$Z_\Phi(s)$.  The second form follows from the zeta functional equation.
\end{proof}

\subsection{Complete large-$y$ expansion}

Let $L=\log2$ and define
\begin{equation}
 C_2=
rac{\zeta''(0)+\tfrac12\log^2(2\pi)-\tfrac{\pi^2}{24}}{L}
     -\frac{L}{12}
 =-1.1090453654341866821907455547\ldots.                  \label{eq:C2}
\end{equation}
For $k\ne0$, let
\begin{equation}
 c_k=-\frac{\Gamma(-\chi_k)\zeta(1-\chi_k)}{L},            \label{eq:ck}
\end{equation}
and define the real analytic, mean-zero, $1$-periodic function
\begin{equation}
 \Psi(u)=\sum_{k\ne0}c_k\e^{2\pi\ii ku}.                  \label{eq:Psi}
\end{equation}

\begin{theorem}[Complex-dimension expansion \NewResult]
\label{thm:H-asymp}
For every $A>0$, as $y\to\infty$,
\begin{equation}
 \boxed{
 H(y)=S-\frac{\log^2S}{2L}-\frac12\log S+C_2
      +\Psi(\log_2S)+\bigO_A(S^{-A}),\qquad S=2\pi y.}    \label{eq:H-asymp}
\end{equation}
The Fourier coefficients decay exponentially.  Numerically,
\begin{align}
 c_1&=(7.5586509397640965-7.4700909117916216\ii)10^{-7},\nonumber\\
 |c_1|&=1.0627109779221651\times10^{-6},\nonumber\\
 |c_2|&=6.6927771176343188\times10^{-13}.                 \label{eq:cnum}
\end{align}
The exact triangle estimate is
\begin{equation}
 \norm{\Psi}_\infty\le2\sum_{k\ge1}|c_k|,
 \qquad
 2\sum_{k\ge1}|c_k|\approx2.1254232944005\times10^{-6}, \label{eq:Psi-bound}
\end{equation}
where the displayed decimal is an 80-digit-working-precision numerical sum.
\end{theorem}

\begin{proofidea}
Shift the contour in \eqref{eq:Mellin-inv} to the left.  The pole at $s=1$
gives $S$.  The triple pole at $s=0$ gives the quadratic logarithm, linear
logarithm, and $C_2$.  The simple poles at $s=\chi_k$, $k\ne0$, give
$c_kS^{\chi_k}=c_k\e^{2\pi\ii k\log_2S}$.  In the functional-equation form,
there are no further poles to the left: the trivial zeta zeros cancel the sine
poles.  Stirling's estimate for $\Gamma(1-s)$ permits a shift to any fixed
negative real part, giving the superalgebraic remainder.  The same estimate
proves exponential Fourier decay.
\end{proofidea}

\subsection{Endpoint duality}

Define the negative Laplace transform of the bounded Fabius variable by
\begin{equation}
 \Lambda(S)=\E\e^{-2SX}.                                   \label{eq:Lambda}
\end{equation}
From $X=(Y+1)/2$ and symmetry of $Y$,
\begin{equation}
 \boxed{\Phi(\ii y)=\e^S\Lambda(S),\qquad S=2\pi y.}       \label{eq:duality}
\end{equation}
Thus \eqref{eq:H-asymp} is exactly the logarithmic large-parameter expansion
which feeds the repository's small-endpoint saddle point.  The Gamma--zeta
coefficients $c_k$ already occur there, but \cref{thm:H-asymp} supplies a new
spectral origin:
\begin{equation}
 \text{endpoint log-periodicity}
 \quad\longleftrightarrow\quad
 \text{residues of }Z_\Phi(s)\text{ at }\chi_k.             \label{eq:duality-slogan}
\end{equation}
The oscillation is tiny not because the dyadic scaling is weak, but because
$\Gamma(-\chi_k)$ is exponentially small at the first nonzero imaginary height.

\section{Finite products, $q$-integers, and pole condensation}
\label{sec:qfinite}

\subsection{Truncated multiplicities}

The $J$-factor product \eqref{eq:PhiJ} has multiplicity
\begin{equation}
 a_J(n)=\min\{J,1+\nu_2(n)\}.                              \label{eq:aJ}
\end{equation}
Let $q=2^{-s}$.  Then
\begin{theorem}[Finite spectral $q$-integer \NewResult]
\label{thm:finite-zeta}
For $\Re s>1$,
\begin{equation}
 Z_{\Phi_J}(s):=\sum_{n\ge1}\frac{a_J(n)}{n^s}
 =\zeta(s)\frac{1-2^{-Js}}{1-2^{-s}}
 =\zeta(s)[J]_q,                                          \label{eq:finite-zeta}
\end{equation}
where $[J]_q=1+q+\cdots+q^{J-1}$.  Also
\begin{equation}
 \sum_{n\le N}a_J(n)=\sum_{h=0}^{J-1}\floor{N/2^h},       \label{eq:AJ}
\end{equation}
while
\begin{equation}
 \log\Phi_J(z)
 =-\sum_{m\ge1}\frac{\zeta(2m)}m[J]_{4^{-m}}z^{2m}
 \qquad(|z|<1).                                           \label{eq:finite-taylor}
\end{equation}
\end{theorem}

Every finite $Z_{\Phi_J}$ is regular at the nonzero $\chi_k$: the apparent
geometric denominator is canceled by the numerator.  The complex dimensions are
therefore genuinely infinite-depth phenomena.

\begin{theorem}[Local pole-condensation profile \NewResult]
\label{thm:condensation}
Fix $k\in\Z$ and put $\chi_k=2\pi\ii k/L$.  Uniformly for $w$ in compact subsets
of $\C$,
\begin{equation}
 \lim_{J\to\infty}\frac1J
 Z_{\Phi_J}\!\left(\chi_k+\frac{w}{J}\right)
 =\zeta(\chi_k)\frac{1-\e^{-Lw}}{Lw},                     \label{eq:condensation}
\end{equation}
with the removable value $1$ for the last factor at $w=0$.
\end{theorem}

\begin{proof}
At $s=\chi_k+w/J$, one has
$q^J=2^{-Js}=\e^{-Lw}$ and
$J(1-q)\to Lw$.  Substitute in \eqref{eq:finite-zeta}.
\end{proof}

The profile $(1-\e^{-Lw})/(Lw)$ is the finite-depth regularization of the limiting
pole.  It gives a precise analytic meaning to the statement that a $q$-integer
turns into a geometric resolvent.

\subsection{Gaussian $q$-binomial layer algebra}

Define the layer zeta functions
\begin{equation}
 \mathcal Z_h(s)=\sum_{2^h\mid n}n^{-s}=q^h\zeta(s).        \label{eq:layer-zeta}
\end{equation}
The finite $q$-binomial theorem gives the exact formal identity
\begin{equation}
 \boxed{
 \prod_{h=0}^{J-1}\bigl(1+t\mathcal Z_h(s)\bigr)
 =\sum_{r=0}^{J}t^r\zeta(s)^r q^{\binom r2}
   \qbinom{J}{r}{q}.}                                     \label{eq:q-layer}
\end{equation}
This is a clean spectral explanation for Gaussian coefficients: they enumerate
independent selections of dyadic zero layers, with weight $q^h$ for choosing
layer $h$.  The identity is deliberately stated as a layer-selection polynomial;
it should not be confused with a Dirichlet series for intersections, because
multiplication of Dirichlet series introduces multiplicative convolution.

\subsection{The common-zero zeta tower}

Intersections are instead encoded by binomial moments of the multiplicity.

\begin{theorem}[Common-zero hierarchy \NewResult]
\label{thm:common-zero}
For every integer $r\ge1$ and $\Re s>1$,
\begin{equation}
 Z_r(s):=\sum_{n\ge1}\binom{a(n)}r n^{-s}
 =\zeta(s)\frac{q^{r-1}}{(1-q)^r},\qquad q=2^{-s}.         \label{eq:Zr}
\end{equation}
For the $J$-layer product,
\begin{equation}
 Z_{J,r}(s)=\sum_{n\ge1}\binom{a_J(n)}r n^{-s}
 =\zeta(s)\sum_{h=r-1}^{J-1}\binom{h}{r-1}q^h.           \label{eq:ZJr}
\end{equation}
Consequently the limiting spectrum has an order-$r$ pole at every nonzero
complex dimension $\chi_k$ unless canceled by $\zeta(\chi_k)$.
\end{theorem}

\begin{proof}
The number $\binom{a(n)}r$ counts $r$-subsets of the layers
$\{0,1,\ldots,\nu_2(n)\}$.  Group such subsets by their largest element $h$.
There are $\binom{h}{r-1}$ choices and their common zeros are precisely the
multiples of $2^h$.  Hence
\[
 Z_r(s)=\zeta(s)\sum_{h\ge r-1}\binom{h}{r-1}q^h
 =\zeta(s)\frac{q^{r-1}}{(1-q)^r}.
\]
Truncating $h$ gives the finite formula.
\end{proof}

This tower is one of the sharpest bridges among the four themes in the user's
request: the sinc layers define the zeros, the $2$-adic intersection counts
produce the rational powers of $(1-q)^{-1}$, and finite layer selections produce
Gaussian $q$-binomials.

\section{The zero heat trace}
\label{sec:heat}

Define the exponentially smoothed zero count
\begin{equation}
 \Theta(t)=\sum_{n\ge1}a(n)\e^{-nt},\qquad t>0.            \label{eq:Theta}
\end{equation}
Layer summation gives the dyadic Lambert series
\begin{equation}
 \Theta(t)=\sum_{h\ge0}\frac1{\e^{2^ht}-1}.               \label{eq:Theta-Lambert}
\end{equation}
Its Mellin representation is
\begin{equation}
 \Theta(t)=\frac1{2\pi\ii}\int_{c-\ii\infty}^{c+\ii\infty}
 \Gamma(s)\frac{\zeta(s)}{1-2^{-s}}t^{-s}\dd s,
 \qquad c>1.                                               \label{eq:Theta-Mellin}
\end{equation}

\begin{theorem}[Heat-trace expansion \NewResult]
\label{thm:heat}
Let $T=\log_2(1/t)$ and
\begin{equation}
 C_{\Theta}=\frac{\gamma-\log(2\pi)}{2\log2}-\frac14
 =-1.1593749760977258236984308878\ldots.                  \label{eq:Ctheta}
\end{equation}
Then, as $t\downarrow0$,
\begin{align}
 \Theta(t)\sim{}&\frac2t-\frac{\log(1/t)}{2\log2}+C_{\Theta}
 +\sum_{k\ne0}\frac{\Gamma(\chi_k)\zeta(\chi_k)}{\log2}
   \e^{2\pi\ii kT}                                      \label{eq:heat-asymp}\\
 &-\sum_{m=0}^{\infty}
 \frac{Z_\Phi(-(2m+1))}{(2m+1)!}t^{2m+1}.\nonumber
\end{align}
The first algebraic terms are
\begin{equation}
 -\frac{t}{12}+\frac{t^3}{5040}+\cdots.                   \label{eq:heat-first}
\end{equation}
The first periodic coefficient has magnitude
$1.2764467470886077\times10^{-6}$.
\end{theorem}

\begin{proofidea}
Shift the contour in \eqref{eq:Theta-Mellin}.  The pole at $s=1$ gives $2/t$.
The double pole at $0$ gives the logarithm and constant.  The poles $\chi_k$
give the periodic term.  Poles of $\Gamma(s)$ at negative even integers are
canceled by trivial zeros of $\zeta$; the negative odd integers give the final
series.
\end{proofidea}

The heat trace makes the spectral dimension picture especially transparent:
linear zero density, a logarithmic digital correction, and a microscopic
log-periodic ripple are successive residue layers of the same Dirichlet series.

\section{A base-$b$ theory}
\label{sec:baseb}

The dyadic case belongs to a larger exact family.  Fix an integer $b\ge2$ and set
\begin{equation}
 \Phi_b(z)=\prod_{h=0}^{\infty}\sinc\!\left(\frac{\pi z}{b^h}\right).       \label{eq:Phib}
\end{equation}
It is the Fourier transform of the $C^\infty$ density obtained by convolving
uniform laws on intervals of half-width $1/(2b^h)$.  Its support is
\begin{equation}
 \left[-\frac{b}{2(b-1)},\frac{b}{2(b-1)}\right].          \label{eq:supportb}
\end{equation}
Let $\nu_b(n)=\max\{h:b^h\mid n\}$, which is meaningful even when $b$ is
composite.

\begin{theorem}[Base-$b$ arithmetic spectrum \NewResult]
\label{thm:baseb}
The zero multiplicity at a nonzero integer is
\begin{equation}
 a_b(n)=1+\nu_b(n),                                       \label{eq:ab}
\end{equation}
and
\begin{align}
 \sum_{n\ge1}\frac{a_b(n)}{n^s}
   &=\frac{\zeta(s)}{1-b^{-s}},                           \label{eq:Db}\\
 \sum_{n\le N}a_b(n)
   &=\sum_{h\ge0}\floor{N/b^h}
     =\frac{bN-s_b(N)}{b-1}.                              \label{eq:Ab}\\
\end{align}
Here $s_b(N)$ is the base-$b$ digit sum.  The entire function has exponential
type $\pi b/(b-1)$ and zero density $b/(b-1)$, in exact agreement with the
Cartwright density law.
\end{theorem}

\begin{corollary}[Even-base automatic lobe signs \NewResult]
If $b$ is even and $x\notin\Z$, then
\begin{equation}
 \sgn\Phi_b(x)=(-1)^{s_b(\floor{|x|})}.                   \label{eq:baseb-sign}
\end{equation}
For arbitrary $b$, the sign sequence
$\sigma_b(N)=(-1)^{\sum_{n\le N}a_b(n)}$ obeys
\begin{equation}
 \sigma_b(bN+r)=(-1)^{bN+r}\sigma_b(N),
 \qquad 0\le r<b.                                         \label{eq:sigmab-rec}
\end{equation}
\end{corollary}

\begin{proof}
Equation \eqref{eq:Ab} has the same proof as the binary digit-sum identity.
For even $b$, its parity equals $s_b(N)$.  The recursion follows directly from
$A_b(bN+r)=bN+r+A_b(N)$.
\end{proof}

\begin{theorem}[Base-$b$ Strang--Fix order \NewResult]
On the lattice $N^{-1}\Z$, the exact Strang--Fix order of the density with
Fourier transform $\Phi_b$ is
\begin{equation}
 1+\nu_b(N).                                               \label{eq:SFb}
\end{equation}
The weighted periodization and Appell reproduction theorems hold verbatim with
$\nu_2$ replaced by $\nu_b$ and the corresponding centered moments.
\end{theorem}

The Mellin calculation also persists.  Let $L_b=\log b$,
$\chi_{b,k}=2\pi\ii k/L_b$, and $S=2\pi y$.  Then
\begin{align}
 \log\Phi_b(\ii y)
 ={}&\frac{b}{2(b-1)}S-\frac{\log^2S}{2L_b}-\frac12\log S+C_b
 +\Psi_b(\log_bS)+\bigO_A(S^{-A}),                        \label{eq:baseb-asymp}\\
 C_b={}&\frac{\zeta''(0)+\tfrac12\log^2(2\pi)-\tfrac{\pi^2}{24}}{L_b}
        -\frac{L_b}{12},                                  \label{eq:Cb}\\
 \Psi_b(u)={}&\sum_{k\ne0}
 \left(-\frac{\Gamma(-\chi_{b,k})\zeta(1-\chi_{b,k})}{L_b}\right)
 \e^{2\pi\ii ku}.                                        \label{eq:Psib}
\end{align}
Thus the Gamma--zeta law is not peculiar to base two; base two is where the lobe
signs become the classical Thue--Morse sequence and the exact dyadic arithmetic
becomes available.

\section{Numerical and symbolic verification}
\label{sec:numerics}

All numerical values in this section were computed independently from the
formulas in the report with 80-decimal working precision.  They are checks, not
proofs.

\subsection{Imaginary-axis residual}

The table compares the direct convergent sum
\begin{equation}
 H(y)=\sum_{h\ge0}\log\frac{\sinh(\pi y/2^h)}{\pi y/2^h}                  \label{eq:H-direct}
\end{equation}
with \eqref{eq:H-asymp}, retaining 20 Fourier modes of $\Psi$.
\begin{center}
\begin{tabular}{@{}r r r@{}}
\toprule
$y$ & $H(y)$ & direct minus asymptotic \\
\midrule
$1$   & $1.8186408379587714$   & $ 3.4873589\times10^{-6}$ \\
$2$   & $5.5707984108206301$   & $ 1.0496\times10^{-11}$ \\
$5$   & $20.0107450110424130$  & $-1.0907\times10^{-12}$ \\
$10$  & $47.2862093875369543$  & $-1.0907\times10^{-12}$ \\
$100$ & $594.0427656248311891$ & $ 2.6185\times10^{-12}$ \\
\bottomrule
\end{tabular}
\end{center}
At $y=1$ the omitted superalgebraic remainder is still visible; by $y=2$ the
error is already near the working evaluation floor.

\subsection{Exact digital checks}

For $N=0,1,\ldots,15$, the cumulative multiplicity
$A(N)=2N-s_2(N)$ has parity
\[
 0,1,1,0,1,0,0,1,1,0,0,1,0,1,1,0,
\]
which is the zero--one Thue--Morse word.  Direct evaluation of 80 sinc factors at
one interior point of each lobe reproduces the corresponding signs.  Exact
integer arithmetic also verifies \eqref{eq:AN}, \eqref{eq:ZJr}, and the
recurrences \eqref{eq:a-rec} through $10^6$ without exception.

\subsection{A reproducible verification recipe}

The main identities can be checked without evaluating $\Up$ itself:
\begin{enumerate}
\item enumerate $a(n)=1+\nu_2(n)$ and compare partial sums with $2N-s_2(N)$;
\item compare $\sum_{n\le M}a(n)n^{-s}$ with
      $\zeta(s)/(1-2^{-s})$ for $\Re s>1$;
\item evaluate $\Phi(x)$ by the sinc product and compare its sign with
      $(-1)^{s_2(\floor x)}$;
\item compare the finite spectral series with $\zeta(s)[J]_{2^{-s}}$;
\item evaluate $H(y)$ both by \eqref{eq:H-direct} and by the residue expansion.
\end{enumerate}
These checks are sufficiently orthogonal that a normalization error in $2\pi$,
a missing initial sinc factor, or an off-by-one valuation is immediately exposed.

\section{What is genuinely added to the pinned corpus}
\label{sec:novelty}

The following table separates established inputs, prior observations, and the
new theorem package.
\begin{longtable}{@{}p{0.37\textwidth}p{0.16\textwidth}p{0.39\textwidth}@{}}
\toprule
Statement & Status & Comment \\
\midrule
\endhead
Canonical product and zero order $1+\nu_2(n)$ & \Corpus & Already in the main
exposition and Lean canonical-product module. \\
All-order endpoint asymptotics and Gamma--zeta coefficients & \Corpus & Already
in the corrected small-argument study. \\
Thue--Morse lobe sign pattern & \Observed & Recorded in the 2018 MSE question;
this report gives the exact zero-count proof and lobe inversion. \\
One critical point per real lobe & \Observed & Essentially proved in the MSE
answer; recast here through Laguerre--Polya theory and the resolvent trace. \\
Arithmetic spectrum package, exact counting law as zero data, spectral traces &
\NewResult & Not found assembled in the pinned TeX corpus. \\
Weighted partition hierarchy and sharp order $1+\nu_2(N)$ & \NewResult & Degree
zero is classical; the exact higher stopping law was not located in the corpus. \\
Appell-corrected exact polynomial reproduction & \NewResult & Uses existing
moment-deconvolution ideas in a different, sharp lattice theorem. \\
Mellin transform of $\log\Phi(\ii y)$ and complex-dimension interpretation &
\NewResult & Explains the already-known endpoint Fourier coefficients by zero-spectrum residues. \\
Finite $q$-integer regularization and pole-condensation profile & \NewResult &
Connects finite sinc products to the limiting imaginary pole lattice. \\
Gaussian $q$-binomial layer-selection polynomial & \NewResult & A spectral-layer
interpretation, distinct from existing dyadic-value $q$-binomials. \\
Common-zero zeta tower $Z_r$ & \NewResult & Encodes binomial multiplicity moments
and higher-order complex-dimension poles. \\
Heat-trace expansion & \NewResult & A residue expansion for exponentially smoothed
zero counts. \\
Base-$b$ extension of the full package & \NewResult & Includes digit-sum signs for
even bases and exact lattice order $1+\nu_b(N)$. \\
\bottomrule
\end{longtable}

\begin{warning}[Historical priority remains narrower than proof]
A targeted search found V.~A.~Rvachev's two-page 1977 paper \emph{On approximation
by means of the function $\Up(x)$}, followed by a substantial atomic-function
approximation literature.  The exact formulas of \cref{thm:weighted,thm:reproduction}
were not located in the pinned corpus or in the limited searchable metadata, but
a Russian-language priority audit is still needed before asserting that those
facts have never appeared elsewhere.
\end{warning}

\section{Lean formalization roadmap}
\label{sec:lean}

The results separate naturally into four dependency layers.

\subsection{Layer I: elementary arithmetic}

Recommended first targets are
\begin{align*}
 a(n)&=\sum_h\mathbf1_{2^h\mid n},\\
 \sum_{n\le N}a(n)&=2N-s_2(N),\\
 (-1)^{a(n)}&=\varepsilon_{n-1}\varepsilon_n,\\
 (-1)^{\sum_{n\le N}a(n)}&=\varepsilon_N.
\end{align*}
The existing Thue--Morse and trailing-zero modules should make these short.
They would yield the lobe-sign theorem once the canonical product's sign changes
are connected to exact orders.

\subsection{Layer II: finite algebra and Dirichlet series}

Finite versions should precede infinite analytic statements:
\begin{itemize}
\item $a_J(n)=\min(J,1+\nu_2(n))$;
\item finite counting formula \eqref{eq:AJ};
\item finite $q$-integer and Gaussian-binomial polynomial identities;
\item finite common-layer count \eqref{eq:ZJr}.
\end{itemize}
For the infinite Dirichlet identities, prove summability on $\Re s>1$ and use
nonnegative layer rearrangement before introducing meromorphic continuation.

\subsection{Layer III: exact lattice reproduction}

The key reusable analytic lemma is a Poisson formula for a compactly supported
$C^\infty$ function and its polynomial multiples.  Then the proof of
\cref{thm:weighted} reduces to:
\begin{enumerate}
\item identify the Fourier coefficients of the periodization;
\item rewrite them as derivatives of $\Phi$;
\item annihilate nonzero coefficients with the existing exact-zero-order theorem;
\item prove failure at the first unannihilated coefficient.
\end{enumerate}
The Appell theorem can be formalized coefficientwise, avoiding formal power-series
asymptotics.

\subsection{Layer IV: complex analysis}

The Mellin and heat-trace theorems need a contour-shift framework with explicit
growth estimates.  A pragmatic sequence is:
\begin{enumerate}
\item formalize the real integral \eqref{eq:kernel-Mellin};
\item prove \eqref{eq:Mellin} only in the fundamental strip;
\item record the residue computations algebraically;
\item defer the all-$A$ remainder until a suitable vertical-contour library is
      available.
\end{enumerate}
Even the strip identity alone gives a valuable formally checked bridge from the
zero spectrum to the endpoint Laplace transform.

\section{Open frontier}
\label{sec:open}

\begin{problem}[Critical-point arithmetic]
Let $c_n$ be the unique critical point in $(n,n+1)$.  Determine asymptotic laws
for $c_n-n$ and $M_n=|\Phi(c_n)|$ conditioned on $\nu_2(n)$, $\nu_2(n+1)$, or
finite binary words around $n$.  The logarithmic derivative
\eqref{eq:logder} converts this into a weighted electrostatic equilibrium problem
with ruler-function charges.
\end{problem}

\begin{problem}[From complex dimensions to real-axis decay]
The imaginary-axis periodic fluctuation is completely controlled by the poles
$\chi_k$, whereas real-axis decay splits into geometric, arithmetic, RMS, and
envelope exponents.  Find an operator or transfer formalism in which both sets of
exponents arise from one spectrum.  A direct residue argument cannot work on the
real axis because zeros and near-zeros dominate.
\end{problem}

\begin{problem}[Sharp approximation constants]
The exact Strang--Fix order gives the first possible nonzero term in a lattice
quasi-interpolant error.  Compute the sharp $L^\infty$, $L^2$, and Sobolev
constants after Appell correction, and determine whether powers-of-two lattices
admit an all-order Richardson scheme whose coefficients have a closed
$q$-binomial form at base $1/4$.
\end{problem}

\begin{problem}[Common-zero zeta residues]
Develop explicit Laurent coefficients of $Z_r(s)$ at $s=\chi_k$.  Their leading
term is elementary, but lower coefficients mix derivatives of zeta with powers
of $\log2$.  Determine whether these coefficients control higher-order
log-periodic corrections of naturally associated spectral statistics.
\end{problem}

\begin{problem}[Base-$b$ exact arithmetic]
For even $b$, the Fourier lobe signs are the generalized digit-sum character
$(-1)^{s_b(n)}$.  Develop the corresponding Fabius-type distribution function,
finite box-spline formulas, exact rational values at $b$-adic points, and
Gaussian-binomial or multinomial expressions.  Composite $b$ should display a
new tension between divisibility by $b^h$ and prime-adic arithmetic.
\end{problem}

\begin{problem}[Inverse spectral characterization]
Classify scale multisets $\{\lambda_h\}$ for which
$\prod_h\sinc(\pi z/\lambda_h)$ has an automatic or regular zero-multiplicity
sequence.  Conversely, characterize which nonnegative automatic sequences can
occur as multiplicities of a locally uniformly convergent geometric sinc
product.
\end{problem}

\section{Conclusion}

The Fabius--Rvachev theory is often presented as a collection of surprising
bridges: a smooth bump from infinite convolutions, rational dyadic values from
$q$-binomials, Thue--Morse signs from finite differences, and log-periodic
asymptotics from a saddle point.  The arithmetic zero spectrum supplies a common
organizing object.

Its summatory law integrates the ruler function into the Thue--Morse lobe signs.
Its exact multiplicities turn Poisson summation into a sharp $2$-adic hierarchy
of polynomial reproduction.  Its Dirichlet series produces the complex
dimensions whose residues are the Gamma--zeta endpoint coefficients.  Its finite
truncations are $q$-integers, its layer selections are Gaussian $q$-binomials,
and its higher intersection statistics are powers of the same geometric
resolvent.  None of these bridges changes the established theory; they explain
why its apparently separate parts keep rediscovering the same dyadic algebra.

The most immediate formal target is \cref{thm:weighted}: it is exact, strong,
short once the current canonical-product theorem is available, and exposes a
clean new role for $\nu_2(N)$.  The most conceptually unifying analytic target is
\cref{thm:H-asymp}, which identifies the endpoint periodic fluctuation as a
residue trace of the Fourier zero spectrum.  Together they turn an infinite sinc
product from a transform formula into an arithmetic spectral object.

\appendix

\section{Residue calculation at the origin}
\label{app:residue}

For completeness, write $L=\log2$ and expand
\begin{align*}
 \zeta(s)&=-\frac12-\frac12\log(2\pi)s+\frac12\zeta''(0)s^2+\bigO(s^3),\\
 \frac1{1-\e^{-Ls}}&=\frac1{Ls}+\frac12+\frac{Ls}{12}+\bigO(s^3),\\
 \frac{\pi}{s\sin(\pi s/2)}&=\frac2{s^2}+\frac{\pi^2}{12}+\bigO(s^2),\\
 y^s&=1+s\log y+\frac{s^2\log^2y}{2}+\bigO(s^3).
\end{align*}
Multiplying and rewriting $\log y=\log S-\log(2\pi)$ gives the residue
\[
 -\frac{\log^2S}{2L}-\frac12\log S+C_2,
\]
with $C_2$ as in \eqref{eq:C2}.  The same calculation for
$\Gamma(s)Z_\Phi(s)t^{-s}$ gives the double-pole contribution
\[
 -\frac{\log(1/t)}{2L}+\frac{\gamma-\log(2\pi)}{2L}-\frac14.
\]

\section{First coefficients and moments}
\label{app:coefficients}

The numerical constants used in the report are
\begin{align*}
 \zeta''(0)&=-2.0063564559085848512101000267299604382\ldots,\\
 C_2&=-1.1090453654341866821907455547094644509\ldots,\\
 C_{\Theta}&=-1.1593749760977258236984308878499458999\ldots.
\end{align*}
For the reciprocal-moment Appell family,
\[
 \frac1{M(z)}=1-\frac{\mu_2}{2!}z^2
 +\frac{6\mu_2^2-\mu_4}{4!}z^4+\cdots,
\]
which with $\mu_2=1/9$ and $\mu_4=19/675$ yields
$6\mu_2^2-\mu_4=31/675$.

\section{Corpus concordance}
\label{app:corpus}

The frozen tree was read recursively.  The most heavily used source families
were the following paths (revision siblings are omitted from this compact list):
\begin{itemize}
\item \nolinkurl{docs/Fabius_Function_and_Rvachev_Up/Fabius_Function_and_Rvachev_Up.tex};
\item \nolinkurl{docs/non-formalized-research-frontiers/non-formalized-research-frontiers.tex};
\item \nolinkurl{docs/archive/standalone-studies/Small_Argument_Asymptotics/Small_Argument_Asymptotics.tex};
\item \nolinkurl{docs/archive/research-frontiers/Fabius_Dyadic_q_Connections/Fabius_Dyadic_q_Connections.tex};
\item \nolinkurl{docs/archive/research-frontiers/Fabius_Dyadic_Formulae_and_Alternative_Representations/Fabius_Dyadic_Formulae_and_Alternative_Representations.tex};
\item \nolinkurl{docs/archive/research-frontiers/Fabius_Thue_Morse_Convergence_Rates/Fabius_Thue_Morse_Convergence_Rates.tex};
\item \nolinkurl{docs/archive/research-frontiers/Repeated_Integration_and_Rvachev_Up/Repeated_Integration_and_Rvachev_Up.tex};
\item \nolinkurl{docs/non-formalized-research-frontiers/drafts/Thue_Morse_Formula_Atlas/Thue_Morse_Formula_Atlas.tex};
\item the complete \nolinkurl{docs/non-formalized-research-frontiers/drafts/rvachev_up_fourier_decay/} revision and audit sequence;
\item primary-paper sources under \nolinkurl{docs/papers/}.
\end{itemize}
The glossary, walkthrough, early notes, primary-exposition revisions, supplements,
$q$-calculus variants, dyadic-asymptotic bridges, inverse-dyadic study, and OCR
siblings were also included in the collision audit.  Their principal role here
was to establish chronology and prevent corrected or archived claims from being
mistaken for independent corroboration.

\begin{thebibliography}{99}

\bibitem{AriasFunction}
J.~Arias de Reyna,
\emph{An infinitely differentiable function with compact support: definition and properties},
arXiv:1702.05442, English translation of the 1982 paper.

\bibitem{AriasArithmetic}
J.~Arias de Reyna,
\emph{Arithmetic of the Fabius function},
arXiv:1702.06487, 2017.

\bibitem{Fabius}
J.~Fabius,
\emph{A probabilistic example of a nowhere analytic $C^\infty$-function},
Z. Wahrscheinlichkeitstheorie verw. Gebiete \textbf{5} (1966), 173--174.

\bibitem{JessenWintner}
B.~Jessen and A.~Wintner,
\emph{Distribution functions and the Riemann zeta function},
Trans. Amer. Math. Soc. \textbf{38} (1935), 48--88.

\bibitem{RvachevApprox}
V.~A.~Rvachev,
\emph{On approximation by means of the function $\Up(x)$},
Dokl. Akad. Nauk SSSR \textbf{233} (1977), no.~2, 295--296;
English translation in Soviet Math. Dokl. \textbf{18} (1977), 340--342.

\bibitem{RvachevSurvey}
V.~A.~Rvachev,
\emph{Compactly supported solutions of functional-differential equations and their applications},
Russian Math. Surveys \textbf{45} (1990), no.~1, 87--120.

\bibitem{ReshetnikovMSE}
V.~Reshetnikov,
\emph{Extrema of an infinite product of sinc functions},
Mathematics StackExchange question 2742261, 2018.

\bibitem{Aistleitner}
C.~Aistleitner, R.~Hofer, and G.~Larcher,
\emph{On evil Kronecker sequences and lacunary trigonometric products},
Ann. Inst. Fourier \textbf{67} (2017), no.~2, 637--687.

\bibitem{AlloucheShallit}
J.-P.~Allouche and J.~Shallit,
\emph{Automatic Sequences: Theory, Applications, Generalizations},
Cambridge University Press, 2003.

\bibitem{StrangFix}
G.~Strang and G.~Fix,
\emph{A Fourier analysis of the finite element variational method},
in \emph{Constructive Aspects of Functional Analysis}, C.I.M.E., 1971, 793--840.

\bibitem{FlajoletMellin}
P.~Flajolet, X.~Gourdon, and P.~Dumas,
\emph{Mellin transforms and asymptotics: harmonic sums},
Theoret. Comput. Sci. \textbf{144} (1995), 3--58.

\bibitem{Levin}
B.~Ya.~Levin,
\emph{Distribution of Zeros of Entire Functions},
American Mathematical Society, revised edition, 1980.

\bibitem{ProveIt}
V.~Reshetnikov,
\emph{ProveIt: Analysis/FabiusFunction}, repository snapshot
\commit, 27 August 2026.

\end{thebibliography}

\end{document}
